% !TEX root = ../../../main.tex
\section{Monotone bounded sequence theorem}
\label{sec:analysis-i:monotone-bounded-sequence-theorem}

% BEGIN TUFTE PLACEHOLDER: supplied sample, lines 548--571.
\begingroup\sloppy
% Replace this entire specimen block with reviewed course notes.
\label{template:sec:page-layout}
\label{template:sec:headings}\index{headings}
This style provides \textsc{a}- and \textsc{b}-heads (that is,
\Verb|\section| and \Verb|\subsection|), demonstrated above.

If you need more than two levels of section headings, you'll have to define
them yourself at the moment; there are no pre-defined styles for anything below
a \Verb|\subsection|.  As Bringhurst points out in \textit{The Elements of
Typographic Style},\cite{Bringhurst2005} you should ``use as many levels of
headings as you need: no more, and no fewer.''

The \TL classes will emit an error if you try to use
\linebreak\Verb|\subsubsection| and smaller headings.

% let's start a new thought -- a new section
\newthought{In his later books},\cite{Tufte2006} Tufte
starts each section with a bit of vertical space, a non-indented paragraph,
and sets the first few words of the sentence in \textsc{small caps}.  To
accomplish this using this style, use the \doccmddef{newthought} command:
\begin{docspec}
  \doccmd{newthought}\{In his later books\}, Tufte starts\ldots
\end{docspec}

\lipsum[4]
\par\endgroup
% END TUFTE PLACEHOLDER
